\documentclass[10pt,a4paper]{article}

% ------------------------------------------------
% Packages
% ------------------------------------------------
\usepackage[utf8]{inputenc}
\usepackage[margin=0.7in]{geometry}

\usepackage{array}
\usepackage{booktabs}
\usepackage{longtable}
\usepackage{multirow}

\usepackage{amsmath}
\usepackage{amsfonts}

\usepackage[table]{xcolor}
\usepackage{colortbl}

\usepackage{titlesec}
\usepackage{microtype}

% ------------------------------------------------
% Color Definitions
% ------------------------------------------------
\definecolor{primary}{RGB}{26,54,93}
\definecolor{secondary}{RGB}{43,108,176}
\definecolor{lightbg}{RGB}{247,250,252}
\definecolor{tablehead}{RGB}{220,230,242}

% ------------------------------------------------
% Typography and Section Styling
% ------------------------------------------------
\titleformat{\section}
{\color{primary}\large\bfseries}
{\thesection}{1em}{}
[\titlerule]

\titleformat{\subsection}
{\color{secondary}\normalsize\bfseries}
{\thesubsection}{1em}{}

% ------------------------------------------------
% Table Formatting
% ------------------------------------------------
\renewcommand{\arraystretch}{1.5}
\setlength{\tabcolsep}{6pt}

% ------------------------------------------------
% Document Title
% ------------------------------------------------
\title{
	\textbf{\color{primary}ER Membrane Protein Classification (Types I--IV)}\\[0.3em]
	\large\textit{
		A Comparative Guide for CSIR NET Life Sciences
	}\\
	\normalsize
	\textit{(Reference: Lodish \& Cooper)}
}

\author{\textbf{Expert High-Yield Reference Summary}}

\date{\small Compiled for CSIR NET Life Sciences Preparation}

% ------------------------------------------------
% Document Begins
% ------------------------------------------------
\begin{document}
	
	\maketitle
	
	% =================================================
	\section{Introduction and Theoretical Framework}
	% =================================================
	
	Integral membrane proteins of the Endoplasmic Reticulum (ER) are classified according to:
	
	\begin{itemize}
		\item Orientation of the N- and C-termini
		\item Presence or absence of cleavable signal peptides
		\item Nature of topogenic sequences
		\item Mechanism of insertion through the \textbf{Sec61 translocon}
	\end{itemize}
	
	For the CSIR NET Life Sciences examination, understanding the following concepts is essential:
	
	\begin{itemize}
		\item \textbf{Positive-Inside Rule}
		\item \textbf{Signal-Anchor (SA) sequences}
		\item \textbf{Stop-Transfer Anchor (STA) sequences}
		\item Membrane topology prediction
	\end{itemize}
	
	% =================================================
	\section{Comprehensive Comparative Table}
	% =================================================
	
	\begin{table}[h!]
		\centering
		\small
		
		\begin{tabular}{|
				>{\raggedright\arraybackslash}m{3.1cm}|
				>{\raggedright\arraybackslash}m{3\documentclass[a4paper,12pt]{article}

\usepackage{ucs}
\usepackage[utf8]{inputenc}
\usepackage{babel}
\usepackage{fontenc}
\usepackage{graphicx}
\usepackage{mhchem}

\usepackage[dvips]{hyperref}

\author{Manpreet Singh}
\date{8/9/26}

\begin{document}

				                                 \end{document}
.1cm}|
				>{\raggedright\arraybackslash}m{3.1cm}|
				>{\raggedright\arraybackslash}m{3.1cm}|
				>{\raggedright\arraybackslash}m{3.1cm}|}
			
			\hline
			
			\rowcolor{tablehead}
			
			\textbf{Feature}
			&
			\textbf{Type I}
			&
			\textbf{Type II}
			&
			\textbf{Type III}
			&
			\textbf{Type IV (Multipass)}
			\\
			
			\hline
			
			\textbf{Pass Status}
			&
			Single-pass
			&
			Single-pass
			&
			Single-pass
			&
			\textbf{Multipass} membrane protein
			\\
			
			\hline
			
			\rowcolor{lightbg}
			
			\textbf{N-terminus Location}
			&
			\textbf{ER lumen}
			&
			\textbf{Cytosol}
			&
			\textbf{ER lumen}
			&
			Variable (IV-A or IV-B)
			\\
			
			\hline
			
			\textbf{C-terminus Location}
			&
			\textbf{Cytosol}
			&
			\textbf{ER lumen}
			&
			\textbf{Cytosol}
			&
			Depends on number of TM segments
			\\
			
			\hline
			
			\rowcolor{lightbg}
			
			\textbf{N-terminal Signal Sequence}
			&
			\textbf{Present and cleavable}
			&
			Absent
			&
			Absent
			&
			Absent
			\\
			
			\hline
			
			\textbf{Topogenic Sequences}
			&
			\textbf{Stop-Transfer Anchor (STA)}
			&
			\textbf{Signal-Anchor (SA)}
			&
			\textbf{Signal-Anchor (SA)}
			&
			Alternating \textbf{SA} and \textbf{STA}
			\\
			
			\hline
			
			\rowcolor{lightbg}
			
			\textbf{Mechanism of Insertion}
			&
			Insertion initiated by N-terminal signal peptide; STA halts translocation and anchors protein
			&
			Internal SA sequence acts as targeting signal and membrane anchor
			&
			Internal SA sequence acts as targeting signal and membrane anchor
			&
			Sequential insertion using alternating SA and STA sequences
			\\
			
			\hline
			
			\textbf{Positive-Inside Rule}
			&
			Positive residues usually remain on cytosolic side
			&
			Positive charges located \textbf{before} SA sequence
			&
			Positive charges located \textbf{after} SA sequence
			&
			Charge distribution determines loop orientation
			\\
			
			\hline
			
			\rowcolor{lightbg}
			
			\textbf{Classic Examples}
			&
			\begin{itemize}
				\item Glycophorin A
				\item LDL receptor
				\item Influenza HA
			\end{itemize}
			&
			\begin{itemize}
				\item Transferrin receptor
				\item Galactosyltransferase
				\item Neuraminidase
			\end{itemize}
			&
			\begin{itemize}
				\item Cytochrome P450
			\end{itemize}
			&
			\begin{itemize}
				\item IV-A: GLUT, Connexin
				\item IV-B: GPCRs (Rhodopsin)
			\end{itemize}
			\\
			
			\hline
			
		\end{tabular}
		
		\caption{
			Structural and mechanistic comparison of Type I--IV ER membrane proteins
		}
		
		\label{tab:er_proteins}
		
	\end{table}
	
	% =================================================
	\section{Core High-Yield Principles for CSIR NET}
	% =================================================
	
	% -------------------------------------------------
	\subsection{1. Positive-Inside Rule}
	% -------------------------------------------------
	
	The orientation of membrane proteins is strongly influenced by the distribution of positively charged amino acids (\textbf{Arginine and Lysine}) surrounding the hydrophobic transmembrane segment.
	
	\begin{itemize}
		\item The side containing more positive charges remains in the \textbf{cytosol}
		\item This rule helps determine final membrane topology
	\end{itemize}
	
	\textbf{Examples:}
	
	\begin{itemize}
		\item \textbf{Type II proteins:}
		Positive charges occur \textbf{before} the SA sequence
		$\Rightarrow$
		N-terminus remains cytosolic
		
		\item \textbf{Type III proteins:}
		Positive charges occur \textbf{after} the SA sequence
		$\Rightarrow$
		C-terminus remains cytosolic
	\end{itemize}
	
	% -------------------------------------------------
	\subsection{2. Topogenic Sequence Definitions}
	% -------------------------------------------------
	
	\begin{itemize}
		
		\item
		\textbf{Signal-Anchor (SA) Sequence:}
		
		An internal hydrophobic sequence that:
		\begin{itemize}
			\item Directs the ribosome to the ER membrane
			\item Interacts with SRP
			\item Remains embedded as a transmembrane helix
			\item Is \textbf{not cleaved}
		\end{itemize}
		
		\item
		\textbf{Stop-Transfer Anchor (STA) Sequence:}
		
		A hydrophobic sequence that:
		\begin{itemize}
			\item Stops translocation through the translocon
			\item Causes lateral release into the lipid bilayer
			\item Forms a membrane-spanning segment
		\end{itemize}
		
	\end{itemize}
	
	% =================================================
	\section{Quick Exam-Oriented Summary}
	% =================================================
	
	\begin{center}
		
		\begin{tabular}{|c|c|c|}
			\hline
			
			\rowcolor{tablehead}
			
			\textbf{Protein Type}
			&
			\textbf{Key Sequence}
			&
			\textbf{Orientation}
			\\
			
			\hline
			
			Type I
			&
			Signal peptide + STA
			&
			N lumen / C cytosol
			\\
			
			\hline
			
			Type II
			&
			Internal SA
			&
			N cytosol / C lumen
			\\
			
			\hline
			
			Type III
			&
			Internal SA
			&
			N lumen / C cytosol
			\\
			
			\hline
			
			Type IV
			&
			Multiple SA + STA
			&
			Multipass topology
			\\
			
			\hline
			
		\end{tabular}
		
	\end{center}
	
\end{document}
